\SourceLeaf{146}
senden, doch erlauben Sie wohl, dass ich es noch einige Wochen behalte. Beim ersten Durchblättern scheint es mir sehr viel Schönes zu enthalten. Vieles von dem, was ich in meiner Methode besonders in ihrer ersten Gestalt an Eigentümlichem hatte, finde ich auch in diesem Buch wieder. Es scheint mein Schicksal zu sein, fast in allen meinen theoretischen Arbeiten mit \textsc{Legendre} zu konkurrieren. So in der höheren Arithmetik, in den Untersuchungen über transzendente Funktionen, die mit der Rektifikation der Ellipse zusammenhängen, bei den ersten Grundlagen der Geometrie und nun wieder hier. So ist z.\,B.\ auch das von mir seit 1794 gebrauchte Prinzip, dass man, um mehrere Größen, die man nicht alle genau darstellen kann, möglichst gut darzustellen, die Summe der Quadrate zu einem Minimum machen muss, auch in \textsc{Legendres} Werk angewandt und recht wacker ausgeführt.

\begin{center}
. . . . . . . . . . . . .
\end{center}

\begin{center}
\textsc{Gauss} an \textsc{Olbers}. Göttingen, 4. Oktober 1809.
\end{center}

. . . . . . . . . Erinnern Sie sich wohl noch, liebster Freund, dass ich bei meinem ersten Aufenthalt in Bremen 1803 mit Ihnen über das Prinzip gesprochen habe, das ich benutzte, um Beobachtungen so genau wie möglich darzustellen, dass nämlich bei gleichem Wert der Beobachtungen die Summe der Quadrate der Differenzen ein Kleinstes sein muss? Dass wir darüber 1804 in Rehburg gesprochen haben, davon sind mir noch alle Umstände gegenwärtig. Es liegt mir daran, dies zu wissen. Über den Grund der Frage ein andermal.

\begin{center}
. . . . . . . . . . . . . . .
\end{center}
